Este proyecto se justifica teniendo en cuenta que esta es una problemática que no se ha visto abordada a profundidad y continua siendo una dificultad para organizaciones de diferentes sectores. Además este proyecto permite aplicar conocimientos de la ingenieria de sistemas como análisis de datos, modelos de predicción y aprendizaje automático para la construcción de una herramienta que contribuya a mejorar los procesos de planificación y administración, generando beneficios tanto para las organizaciones y los usuarios finales.

\subsubsection{Impactos}

Mejorar la disponibilidad de los productos mediante una mejor planificación de los inventarios, apoyar la toma de decisiones del personal encargado del abastecimiento,y fortalecer el uso de herramientas tecnológicas en el sector administrativo.
